\par\bigskip\SourceLeaf{163}
\Needspace{8\baselineskip}\pdfbookmark[1]{III. Commission report}{witwen-III}
\begin{center}
[III.]\par\medskip
[\textit{Commission report.}]
\end{center}
\noindent\hspace*{\parindent}The commission appointed to consider the affairs of the widows' fund has the honour to present the outcome of its deliberations.

The first question was what rate of interest should be taken as the basis for the measures to be adopted.

Most of the capital of the widows' fund is at present still invested at a rate above \(3\tfrac{1}{2}\) per cent; nevertheless, a proper assessment of the circumstances requires the average rate of interest to be regarded as definitely below 4 per cent. But, having considered the following grounds:

1) that the rate of interest in every country of Europe shows an unmistakable tendency to fall further, despite individual fluctuations\footnote{One of the most enlightened statesmen, especially in financial matters, the Grand Ducal Baden minister \textsc{Nebenius}, says in his well-known work Über die Herabsetzung der Zinsen der öffentlichen Schulden, p. 129: ‘The gradual further fall in the rate of interest, which cannot fail to occur if peace continues for a longer time, will eventually bring about a reduction to three per cent everywhere, here a little earlier, there a little later, or at least make it appear a measure depending only on the government's decision’; and elsewhere, p. 21: ‘A period of greater activity in productive undertakings, which checks the fall in the rate of interest for a time, is followed all the more certainly by a rapid fall in the rate of interest.’};

2) that every institution which depends essentially on the rate of interest must, if it is not to be regarded as entirely insecure, be based on a rate somewhat below the one prevailing at the moment;

3) that the regulations of 1794 also set us a good example, inasmuch as there was at least an intention at that time to make the fund secure at an interest rate of 3 per cent;

\noindent the commission considers it inadmissible to take a rate above \(3\tfrac{1}{2}\) per cent as the basis, and therefore considers it absolutely necessary to aim to cover the resulting deficit of 17520 thalers at this rate.

If it were thought that the annual surpluses still arising at present under the existing arrangements {\renewcommand{\thefootnote}{**)}\footnote{The annual account for 1844—1845 shows no surplus, since
\par\noindent for 1844 July 1 the monetary assets were stated as 116375 Thl. 14 Ggr. 6 Pf.
\par\noindent but for 1845 July 1, only as . . . . . . . . . . . . . . . . . 116369 ,, 17 ,, 4 ,,
\par\noindent Since, however, there are three fewer widows in the current year, a surplus is again to be expected this year, and will probably continue for some time.}} were already in themselves steps towards covering the deficit, and that it would have been extinguished as soon as the present assets had increased by 17520 thalers, such an opinion would rest only on a confusion of concepts and a complete misunderstanding of the meaning of the deficit. The result of the calculation by which the size of the deficit was determined already implicitly takes full account of the annual savings still occurring at present; to regard these as steps towards covering it would be the same as counting them twice {\renewcommand{\thefootnote}{***)}\footnote{Since similar views have reportedly been expressed here and there, but this report would not be the place to refute them, a further discussion of the subject has been added to the appendix, although it will be entirely superfluous for all who have already read and considered the memorandum with the necessary attention.}}. Let the association therefore entertain no illusion about the truth
\SourceLeaf{164}
that the deficit can and must be covered solely by suitable changes to the existing arrangements.

It is self-evident that weak measures can produce only weak effects. These would include making participation, hitherto voluntary, compulsory for all professors appointed in future, and removing the freedom to withdraw again at any time. That such a measure could have only a very slight effect is evident from the fact that, over a 50-year average, the average number of professors who did not contribute although entitled to participate was only \(=4{,}4\); hence the annual income thereby lost to the fund at the existing level of contributions was \(=44\) thalers, which accordingly represents permanent capital of 1100 or 1257 Rthl., depending on the rate of interest. But this is offset, if not wholly, at least in large part, by the penalties payable for delayed entry under the existing, very suitable arrangement. Some of these have been very substantial, as the following examples show. The following have made back payments at double the rate:
\begin{center}
L. for 12 years,\quad R. for 13 years,\quad L. for 19 years,\quad S. for 25 years,\quad \mbox{C. for 30 years.}
\end{center}
\noindent\hspace*{\parindent}The commission therefore considers that so small and doubtful an advantage as might accrue to the fund from turning the institution into a compulsory one cannot weigh against the destruction of the foundation's hitherto liberal character.

A couple of other measures whose effect is likewise only weak or uncertain will be mentioned at the end of this report.

Accordingly, only the following can be regarded as truly effective measures:

1) Increasing the contributions.

2) Reducing the pensions.

3) Combining both measures.

The balance calculation set out in the memorandum makes it possible to state exactly to what extent these measures must be applied if the aim is to be achieved, that is, if the balance is to show an improvement of 17520 Rthl. Details would be out of place here; the final result, however, is:

1) If the deficit is to be covered solely by increasing the contributions, these must be raised universally, that is, for \textit{all} present and future members, to \(4\tfrac{1}{4}\) Louis d'or.

2) If the reduction of pensions alone is to cover it, they must be reduced to 223 Rthl. (while that of Professor G.'s widow would remain at 200 Rthl.).

3) If the burden is shared between contributors and pensioners, the result would depend on the chosen proportion: if, for example, the former were to bear \(\tfrac{2}{3}\) and the latter \(\tfrac{1}{3}\), the contributions would have to be \(3\tfrac{1}{2}\) Louis d'or and the pensions 241 Rthl. If the contributions are to be raised only to 3 Louis d'or, then, if the deficit is really to be covered, only a pension of 235 Rthl. can be paid.

All these calculations assume that the chosen changes take effect from Michaelmas 1845, and that, in the second or third case, the reductions in pensions apply to all widows, present and future alike. If, without any increase in contributions and without a reduction in pensions for the present widows, the future widows were to bear the burden alone, only a pension of \(213\tfrac{1}{4}\) Rthl. could be granted to them.

